%% shortname: State-Contingent Wealth, Insurance and First-Order Stochastic Dominance
\documentclass[11pt]{article}
\usepackage{microstyle}

\begin{document}

\lessonheader{18}{Chapter 6 (continued): state-contingent wealth, insurance, risky assets and first-order stochastic dominance}{}


\section{Where we are}

We started looking at lotteries over wealth. The properties of the Bernoulli utility
function $u(\cdot)$ tell us whether the agent is risk averse or not.

Our aim is to evaluate and compare different lotteries \emph{without} assuming too much
knowledge about the Bernoulli utility function.


\section{A different interpretation of the model: state-contingent wealth}

But first, a different interpretation of the model. Rather than changing the
\emph{probabilities} in lotteries, we change the \emph{wealth level} in different
states (outcomes); e.g.\ insurance, risky assets.

\topic{Setup}
With a finite number of states and a fixed probability distribution
$p=(p_1,\dots,p_N)$, the expected utility of the wealth vector $x=(x_1,\dots,x_N)$ is
\[ V(x)=p_1u(x_1)+\dots+p_Nu(x_N)=\sum_{n=1}^{N}p_n u(x_n). \]

\topic{Risk aversion $\Rightarrow$ convex preferences}
If the agent is risk averse, $u(\cdot)$ is concave, and $V$ is a sum of concave
functions with non-negative weights $p_n$. Hence
\[ V(x)\ \text{is concave}\ \Rightarrow\ V(x)\ \text{is quasiconcave}\
   \Rightarrow\ \text{preferences over } x \text{ are convex.} \]
$\Rightarrow$ Averages are preferred to extremes. This fits with risk aversion: we like
more balanced $x$'s because they have less risk, i.e.\ less variability across
$x_1,\dots,x_N$.

\subsection*{An example with 2 states}

State $1$ is the good state, state $2$ the bad state.

\begin{center}
\begin{tikzpicture}[scale=1.0]
  \ecaxes{5.2}{5.2}
  \node[ecnode,right] at (5.2,0) {$x_1$, good state};
  \node[ecnode,above] at (0,5.2) {$x_2$, bad state};
  \draw[dashed] (0,0) -- (4.8,4.8);
  \node[ecnode,anchor=west,align=left] at (4.85,4.8) {line of certainty\\(full insurance)};
  \node[ecnode] at (0.75,0.35) {$45^\circ$};
  % indifference curve x1 x2 = 4
  \draw[budget,domain=0.85:4.7,samples=60] plot (\x,{4/\x});
  \node[bundle] at (0.85,4.706) {}; \node[bundle] at (4.7,0.851) {};
  \node[bundle] at (2,2) {};
  \draw[thin] (0.8,3.2) -- (3.2,0.8);
  \node[ecnode,anchor=north west] at (3.25,0.75) {slope $-\tfrac{p_1}{p_2}$};
  \node[ecnode,right] at (4.7,0.95) {$V(x)=\bar V$};
  \draw[-{Stealth}] (2.9,1.38) -- (3.4,1.88);
  \node[ecnode,right] at (3.4,1.95) {preferred};
\end{tikzpicture}
\end{center}

\[ V(x)=p_1u(x_1)+p_2u(x_2), \qquad
   \mathrm{MRS}=-\frac{p_1u'(x_1)}{p_2u'(x_2)} \quad\text{(slope of the indifference curve).} \]

\begin{flatlist}
  \item If the agent is risk neutral, $u(x)=x$, so
        $\mathrm{MRS}=-\dfrac{p_1}{p_2}$, a constant: indifference curves are straight
        lines.
  \item For a risk-averse agent, \emph{on the $45^\circ$ line} ($x_1=x_2$, so
        $u'(x_1)=u'(x_2)$):
        \[ \mathrm{MRS}=-\frac{p_1}{p_2}. \]
\end{flatlist}
Starting at a point of certainty, $x_1=x_2$, it is as if the agent were risk neutral
for small changes.


\section{Application to insurance}

The price of $1$ unit of insurance is $\alpha$. If you buy $1$ unit, you pay $\alpha$ in
both states and receive $1$ in the bad state:
\[ dx_1=-\alpha, \qquad dx_2=1-\alpha, \qquad
   \frac{dx_2}{dx_1}=\frac{1-\alpha}{-\alpha}=-\frac{1-\alpha}{\alpha}. \]

\topic{Overpriced insurance}
For the insurance company to make a profit we need $\alpha\ge p_2$, so
\[ \frac{1-\alpha}{\alpha}\ \le\ \frac{1-p_2}{p_2}=\frac{p_1}{p_2}: \]
the budget line is flat compared to the indifference curve at the $45^\circ$ line.

\begin{center}
\begin{tikzpicture}[scale=1.0]
  % p1 = p2 = 1/2, u = ln; budget x2 = 3.32 - 0.6 x1 through e = (4.2, 0.8)
  \fill[black!8] (0,0) -- (0,3.32) -- (5.533,0) -- cycle;
  \ecaxes{6.0}{4.6}
  \node[ecnode,right] at (6.0,0) {$x_1$};
  \node[ecnode,above] at (0,4.6) {$x_2$};
  \draw[dashed] (0,0) -- (4.3,4.3);
  \node[ecnode] at (0.75,0.35) {$45^\circ$};
  \draw[budget] (0,3.32) -- (5.533,0);
  \node[ecnode,anchor=west] at (5.2,0.75) {slope $-\tfrac{1-\alpha}{\alpha}$};
  % IC through the certain point on the budget line (x1 x2 = 2.075^2)
  \draw[domain=1.05:4.2,samples=60] plot (\x,{4.3056/\x});
  \draw[thin] (1.2,2.95) -- (2.95,1.2);
  \node[ecnode,anchor=north east] at (2.95,1.2) {$-\tfrac{p_1}{p_2}$};
  \node[bundle] at (2.075,2.075) {};
  % optimal IC x1 x2 = 4.593, tangent at (2.767, 1.66)
  \draw[budget,domain=1.15:5.3,samples=60] plot (\x,{4.593/\x});
  \node[bundle] at (2.767,1.66) {};
  \node[ecnode,above right] at (2.767,1.66) {$x^*$};
  \node[bundle] at (4.2,0.8) {};
  \node[ecnode,below left] at (4.2,0.8) {$e$};
  \node[ecnode] at (1.75,0.35) {feasible set};
\end{tikzpicture}
\end{center}

$\Rightarrow$ Even a risk-averse agent accepts some risk (does not fully insure) if
insurance is overpriced: the optimum $x^*$ lies below the $45^\circ$ line, $x_1^*>x_2^*$.
(Here $e$ is wealth with no insurance.)


\section{Application: betting on a risky asset}

For each unit of the asset bought,
\[ \left.\begin{aligned} dx_1&=\alpha_1>0\\ dx_2&=-\alpha_2<0 \end{aligned}\right\}
   \qquad \frac{dx_2}{dx_1}=-\frac{\alpha_2}{\alpha_1}. \]
The expected return is
\[ p_1\alpha_1+p_2(-\alpha_2), \]
and the return is positive iff
\[ p_1\alpha_1>p_2\alpha_2 \quad\Longleftrightarrow\quad \frac{p_1}{p_2}>\frac{\alpha_2}{\alpha_1}. \]
Then, starting from certainty, the indifference curve (slope $-p_1/p_2$) is steeper than
the asset line (slope $-\alpha_2/\alpha_1$).

$\Rightarrow$ The risk-averse agent buys a small (but positive) amount of the risky
asset whenever its expected return is positive.


\section{Comparing distributions (lotteries)}

If we know the agent's Bernoulli utility function, it is easy to pick between two
lotteries: pick the one with the highest $U(F)$, or equivalently the highest CE.

But what if $u(x)$ is not known? We look for \emph{robust} comparisons, where it does
not matter what the exact details of $u(x)$ are.

$\Rightarrow$ Are there lotteries that ``everyone'' can agree are better than others?

\runin{Aside.} In Micro II, we use distributions to describe beliefs about
characteristics (ability, willingness-to-pay). Being able to compare different
distributions has many applications.

\topic{Setting}
We compare distributions with the same support $[0,\bar x]$, i.e.\ with
\[ F(0)=0, \qquad F(\bar x)=1. \]

\runin{First question.} Assume the only information we have is that more is better
than less, i.e.\ $u(x)$ is non-decreasing. Are there distributions $F$ and $G$ such
that all these agents prefer $F$ to $G$?


\section{First-order stochastic dominance (FOSD)}

\topic{Definition 1}
$F$ first-order stochastically dominates $G$ if
\[ \int u(x)\,dF(x)\ \ge\ \int u(x)\,dG(x) \qquad\text{for all non-decreasing } u(x).
   \tag{1} \]

\topic{Definition 2}
$F$ first-order stochastically dominates $G$ if
\[ F(x)\ \le\ G(x) \qquad \forall\, x\in[0,\bar x]. \tag{2} \]

\topic{Main result}
$(1)\Longleftrightarrow(2)$: the two definitions are equivalent.

\runin{Intuition for $(2)\Rightarrow(1)$.} $F(x)\le G(x)$ means that there is a smaller
probability of a bad outcome (at most $x$) under $F$ than under $G$
\[ \Longleftrightarrow\quad 1-F(x)\ \ge\ 1-G(x): \]
$F$ is more likely to give a good outcome (more than $x$).

\begin{center}
\begin{tikzpicture}[scale=1.45]
  % G(x) = 1-(1-x/4)^2, F(x) = (x/4)^2, height 1 = 2.6
  \ecaxes{4.8}{3.2}
  \node[ecnode,right] at (4.8,0) {$x$};
  \node[ecnode,below left] at (0,0) {$0$};
  \draw[dashed] (0,2.6) -- (4,2.6);
  \node[ecnode,left] at (0,2.6) {$1$};
  \draw[budget,domain=0:4,samples=60] plot (\x,{2.6*(1-(1-\x/4)^2)});
  \draw[budget,domain=0:4,samples=60] plot (\x,{2.6*(\x/4)^2});
  \node[ecnode,anchor=north west] at (2.6,2.25) {$G(x)$};
  \node[ecnode,anchor=north west] at (3.0,1.35) {$F(x)$};
  \draw[dashed] (4,0) node[ecnode,below] {$\bar x$} -- (4,2.6);
  \draw[dashed] (1.5,0) node[ecnode,below] {$x$} -- (1.5,2.6);
  \draw[{Stealth}-{Stealth}] (1.35,1.584) -- (1.35,2.6);
  \node[ecnode,left] at (1.35,2.1) {$1-G(x)$};
  \draw[{Stealth}-{Stealth}] (1.65,0.366) -- (1.65,2.6);
  \node[ecnode,right] at (1.65,1.25) {$1-F(x)$};
\end{tikzpicture}
\end{center}

\topic{Proof of $(2)\Rightarrow(1)$}
Assume $u(x)$ is differentiable, and compare expected utilities:
\[ \int u(x)\,dF(x)-\int u(x)\,dG(x)=\int u(x)\,dH(x), \qquad
   H(x)=F(x)-G(x)\ \le\ 0. \]
Integrating by parts,
\begin{align*}
  \int_0^{\bar x} u(x)\,dH(x)
   &= \bigl[u(x)H(x)\bigr]_0^{\bar x}-\int_0^{\bar x} u'(x)H(x)\,dx\\
   &= 0-0-\int_0^{\bar x} \underbrace{u'(x)}_{\ge0}\,\underbrace{H(x)}_{\le0}\,dx
    \ \ge\ 0 \qquad\text{if } u'(x)\ge 0,
\end{align*}
since $H(0)=F(0)-G(0)=0$ and $H(\bar x)=1-1=0$. So (1) holds. $\checkmark$

\topic{Proof of $(1)\Rightarrow(2)$}
$F$ is better than $G$ for \emph{all} non-decreasing utility functions, and therefore
also for the step function
\[ u(x)=\begin{cases} 0 & \text{if } x\le z,\\ 1 & \text{if } x>z. \end{cases} \]
From (1),
\begin{align*}
  0\ \le\ \int u(x)\,dH(x) &= \int_0^{z} 0\,dH(x)+\int_z^{\bar x} 1\,dH(x)
   = \bigl[H(x)\bigr]_z^{\bar x}\\
  &= 0-\bigl(F(z)-G(z)\bigr)=G(z)-F(z).
\end{align*}
And $z$ does not matter $\Rightarrow$ $F(z)\le G(z)$ for all $z\in[0,\bar x]$
$\Rightarrow$ (2). $\checkmark$


\begin{transcriptnotes}
  \item The export has 7 pages; page 7 holds only the last two lines of the proof
        (the rest of the page is blank). No date or title is written, so the header
        date is left empty; the subtitle and section headings are \textbf{added}
        (MWG Chapter 6, Sections 6.C--6.D).
  \item Page 1: ``benoulli'' $\to$ \emph{Bernoulli}; ``evaluate compact differend
        lotteries'' $\to$ \emph{evaluate and compare different lotteries};
        ``differend'' $\to$ \emph{different}; ``probilities'' $\to$
        \emph{probabilities}; ``funtions'' $\to$ \emph{functions}. ``less variability
        because of $x_1\dots x_n$'' $\to$ \emph{across $x_1,\dots,x_N$}. The page
        writes $V(x)=p_1u(x_1)\ \dots\ p_Nu(x_N)$ with the $+$ signs missing; the sum
        is restored. ``with non-negative weights $p_n$'' is \textbf{added} (a sum of
        concave functions is concave only with non-negative weights).
  \item Page 2 diagram (two states, $45^\circ$ line labelled ``line of uncertainty /
        full insurance''): \textbf{corrected} to \emph{line of certainty} --- on the
        $45^\circ$ line wealth is the same in both states. Redrawn to scale with
        $p_1=p_2$ and $u=\ln$; the tangent with slope $-p_1/p_2$ on the $45^\circ$
        line and the ``preferred'' arrow (the page has small arrows) are drawn
        explicitly. The page writes $U(x)$ here and $V(x)$ on page 1; typeset $V(x)$
        throughout.
  \item ``Straight line: for a risk averse agent, on the $45^\circ$ line
        $\mathrm{MRS}=-p_1/p_2$'' --- the ``straight line'' belongs to the
        risk-neutral case; split into two bullets. ``for some changes'' $\to$
        \emph{for small changes} (the statement is local).
  \item Insurance: the top of page 3 is cut off; it shows ``$\ldots=-\alpha$,
        $dx_2=1-\alpha$'', read as $dx_1=-\alpha$, $dx_2=1-\alpha$. The explanation
        ``pay $\alpha$ in both states and receive $1$ in the bad state'' is
        \textbf{added}. The step $(1-p_2)/p_2=p_1/p_2$ uses $p_1+p_2=1$.
  \item Insurance diagram redrawn to scale ($p_1=p_2=\tfrac12$, $u=\ln$,
        $\alpha=0.625$, so slope $-0.6$). The no-insurance point $e$ and the label
        $x^*$ are \textbf{added}. ``accents risks'' $\to$ \emph{accepts (some)
        risk}; ``does not fully insure'' and $x_1^*>x_2^*$ are \textbf{added}.
  \item Risky asset: ``Application hetting a risk asset'' $\to$ \emph{betting on a
        risky asset}; ``asset but'' $\to$ \emph{asset bought}; ``expected winning of
        return'' $\to$ \emph{expected return}. ``buys a small amount of risk
        asset'' $\to$ \emph{a small (but positive) amount \ldots\ whenever its
        expected return is positive}; the steeper-indifference-curve sentence is
        \textbf{added} to connect it to the $45^\circ$-line MRS.
  \item Page 4: ``companies distribution lotteries'' $\to$ \emph{Comparing
        distributions (lotteries)}. ``highest $U(F)$ on CE'' $\to$ \emph{or the
        highest CE}. ``cability'' $\to$ \emph{ability}; ``into'' $\to$
        \emph{information}; ``all these agent $F$ to $G$'' $\to$ \emph{all these
        agents prefer $F$ to $G$}.
  \item Page 5: ``FOSD'' expanded to \emph{first-order stochastically dominates}. In
        Definition 1 the right-hand side is written $\int u\,dG(x)$; typeset
        $\int u(x)\,dG(x)$. ``small prob of a bad under $F$'' $\to$ \emph{smaller
        probability of a bad outcome (at most $x$)}; the line
        $1-F(x)\ge1-G(x)$ is \textbf{added}. The CDF sketch (with ``$1-G(x)$''
        marked) is redrawn to scale with $G(x)=1-(1-x/\bar x)^2$,
        $F(x)=(x/\bar x)^2$; the $1-F(x)$ gap is \textbf{added}.
  \item Page 6, $(2)\Rightarrow(1)$: ``differentio'' $\to$ \emph{differentiable}.
        The boundary terms vanish because $H(0)=H(\bar x)=0$; this reason is
        \textbf{added}.
  \item Page 6, $(1)\Rightarrow(2)$: the step function is named $v(x)$ on the page but
        then used as $u(x)$ in the next line; typeset $u(x)$ throughout. Note that
        it is not continuous, whereas Lesson 16 assumed a continuous $u$; the proof
        uses Definition 1 for all non-decreasing $u$ (a continuous approximation
        gives the same conclusion). $[H(x)]_z^{\bar x}=H(\bar x)-H(z)$ is the
        integral over $(z,\bar x]$, consistent with the cut at $x\le z$.
\end{transcriptnotes}

\end{document}
